%%%%%%%%%%%%%%%%%%%%%%%%%%%%%%%%%%%%%%%%%%%%%%%%%%%%%%%%%%%%%%%%%%%%%%%%%%
%
% Template for making new Astro Inferno documents with the Infernal LaTeX
%
%%%%%%%%%%%%%%%%%%%%%%%%%%%%%%%%%%%%%%%%%%%%%%%%%%%%%%%%%%%%%%%%%%%%%%%%%%

%  * * * Remember to comply with the Astro Inferno 3rd party licence! * * *
% Somewhere in the document body you should put an Astro Coven logo
% and the 3rd party licence text.
% The former is included with the template
% (img/astro-inferno-logo-kit/Astro~Inferno~-~Coven~[Black|White].png).
% \AIcopyright or \AIcopyright* wiill print the licence text (starred version circumvents
% problems with missing copyright symbols in the font used to replace 1651 Alchemy).
%
% The pre-made titlepage and table of contents (\AIMakeElementTitlepage, \AIMakeElementTOCpage)
% renders both 3rd party licence text and coven logo.
%
% Acknowledging the Infernal LaTeX is not required, but appreciated!
%
% It is your responsibility to make sure you comply with the font licences.
% And of course that you have the rights to use any content you put into this template.
%  * * *

%%%%%%%%%%%%%%%%%%%%%%%%%%%%%%%%%%%%%%%%%%%%%%%%%%%%%%%%%%%%%%%%%%%%%%%%%%
%
%   ===~- Preamble -~===
%

%  * * * Here we invoke the astroinfernobook style. * * *
%
% The provided options mimic Astro Inferno.
%
% Change option letterpaper to a4paper to get A4 paper size better for printing in Europe)
%
% Option twocolumn is strongly recommended. Astro Inferno uses two columns most of the time
% and the template was designed with that in mind. Some elements might not be tested
% in onecolumn mode.
%
\documentclass[letterpaper,twocolumn,openany,nodeprecatedcode]{astroinfernobook}

% Load some packages. Add more if needed.
\usepackage[utf8]{inputenc}
\usepackage[singlelinecheck=false]{caption}
\usepackage{stfloats}
\usepackage[hidelinks]{hyperref}

% Control table captions
\captionsetup[table]{labelformat=empty,font={sf,sc,bf,},skip=0pt}

% Load Berylium font
% \newfontfamily\AIfontBeryllium{Berylium_Rg.otf}[
%     Path = fonts/berylium/,
%     Extension = .otf,
%     UprightFont = *,
%     BoldFont = Berylium_Bd,
%     ItalicFont = Berylium_Rg_It,
%     BoldItalicFont = Berylium_Bd_It
% ]

%
%
%%%%%%%%%%%%%%%%%%%%%%%%%%%%%%%%%%%%%%%%%%%%%%%%%%%%%%%%%%%%%%%%%%%%%%%%%%
%
%   ===~- The actual body of the document -~===
%

%  * * * Dark and Light mode * * *
% The document starts in light mode by default. If you want to start in dark mode, comment out \AIMakeDark* below.
% \AIMakeDark*
%
% To switch mode later in the document, use \AIMakeDark and \AIMakeLight
% where you want to make a page break and continue in the new mode.


\begin{document}

% Start frontmatter, which is titlepage, TOC, etc.
% Uses unnumbered chapters, roman numerals for page numbers.
\frontmatter


% You want a title page?
\begin{titlepage}
  % You can make your own title page however you like.
  % Or you can use the pre-made one, which is a bit flexible via options etc.
  % The pre-made one prints the mandatory 3rd party licence text.
  \AIMakeElementTitlepage
  [
  % You can use titlefont to pick a font for the title block. AvalenRekas is the default.
  % titlefont = \AIfontBeryllium,
  quote = {\AIQuote[width=0.6\textwidth, source=- Unholy Diver, sourcegap=\baselineskip]
    {
      This is where you can put a verse or quote.\\
      See {\normalfont \texttt{lib/astroinfernoutility.sty}} for options and default values.\\
      Left-adjusted quote, right adjusted source by default.
    }
  }
  ]
  {% Ttile block.
    {\Huge Title\\[0.5em]}
    % Using the default AvalenRekas for the smaller stuff becomes a bit heavy, so we switch to normalfont (lato)
    \normalfont
    {\large Subtitle\\[1.5em]}
    Author name\\
    \today
  }
\end{titlepage}


% You want a table of contents page?
% You get the pre-made one with \AIMakeElementTOCpage, including mandatory coven logo.
% Just a normal \tableofcontents actually goes a long way too.
\AIMakeElementTOCpage


% Time for the real content. mainmatter starts numbering chapters, restarts page numbering using arabic numerals.
\mainmatter


% Start new part
\part{Commence brutality}

% Before starting a new chapter, we add inject some empty space above and a padded word of warning below the chapter title.
% Here, we use \AIQuote for the latter, which is not needed, you can just enter text or other things if you want.
\AINextChapterBefore[2cm]{}
\AINextChapterAfter[1cm][1cm]{
  \AIQuote[adjust=c]{
    \textbf{Warning!}\\[\baselineskip] % Add extra space between lines
    Chapter might contain exciting material.\\
    Not suitable for dull people.
  }
}
% Actually start the chapter. Give it a label if you want to refer to it elsewhere.
\chapter{Soul evisceration} \label{ch:slaughter}

% Start a new section
\section{Materials} \label{sec:materials}
\begin{itemize}
\item 1 Soul
\item Spirit cutter
\item Sturdy table
\item Shackles or restraints
\item Coal
\end{itemize}
\section{Preparations}
% Start writing away
First, bring out a newly sharpened spirit cutter (see \nameref{sec:materials}). Heat over burning coal until red-hot.

\ldots

\end{document}

%
% Once finished editing, compile with lualatex (here assuming UNIX-like system):
% lualatex astrobase.tex && lualatex astrobase.tex
%
% See how we ran it twice through lualatex, which is needed to get references right, also borders.
%
% Alternatively use latexmk. If working in the directory of the Infernal LaTeX, the Makefile makes this simple:
% make astrobase.pdf